\documentclass[10pt]{report}
\usepackage{geometry}
\geometry{a4paper, top=20mm, left=30mm, right=25mm, bottom=30mm, headsep=10mm, footskip=12mm}
\usepackage[utf8]{inputenc}

%Options: Sonny, Lenny, Glenn, Conny, Rejne, Bjarne, Bjornstrup
\usepackage[Bjornstrup]{fncychap}
\usepackage{ngerman}
\usepackage{graphicx}

\begin{document}
\title{Pflichtenheft}
\begin{figure}
	\begin{flushright}
		\includegraphics[scale=0.5]{uniLogo.png}
	\end{flushright}
\end{figure}
\begin{center}
\vspace{2.0 cm}
{\LARGE SEP – Wintersemester 2013/14}

\vspace{1.0 cm}
\textbf{{\Huge Pflichtenheft}}

\vspace{0.5 cm}
\begin{figure}[!htb]
\begin{center}
	\includegraphics[scale=1.0]{projektLogo.eps}
\end{center}
\end{figure}

\vspace{0.5 cm}
\textbf{{\Huge OpenStreetMap: Die Welt in 3D}}

\vspace{1.5 cm}
[endgültiges Datum]

\vspace{0.5 cm}
Version: 1.0

\vspace{1.5 cm}
{\Large Projektbetreuer: Peter Barth}

\vspace{1.5 cm}
\begin{tabular}{c c c}
Phase & Verantwortlicher & E-Mail Adresse \\
\hline \hline
Pflichtenheft & Haas Gabriele & haasgab@fim.uni-passau.de \\
\hline \hline
Entwurf & Eder Thomas & \\
\hline \hline
Spezifikation & Blauberger Christof & blauberg@fim.uni-passau.de \\
\hline \hline
Implementation & Knorr Fabian & \\
\hline \hline
Testing & Wenger Constantin & \\
\hline \hline
Präsentation & Reichl Sebastian & \\
\end{tabular}


\end{center}
Daten: Version Nr. (Datum + Uhrzeit), erstellt von: Team-Mitglieder

\tableofcontents

nicht tiefer als 3 - 4 Gliederungsebenen

\chapter{Ausgangssituation}
[Wie stellt sich die Situation heute dar?]

Die Kartendaten des OpenStreetMap-Projekts erfreuen sich immer größerer Beliebtheit. Der Detailgrad der Daten, die Menge an verschiedenen Merkmalen und die Genauigkeit der Daten sind in den meisten Regionen der Welt ihrer Konkurrenz weit voraus. Durch die verschiedenartigen Karten, Kartenstile und Spezialanwendungen auf Basis von OpenStreetMap gibt es eine unzählige Menge von Anwendungsmöglichkeiten.


[Was ist der Auslöser für die Erstellung des Pflichtenhefts?]

Aus diesem Grund soll eine 3D-Desktopanwendung basierend auf den Daten von OpenStreetMap erstellt / implementiert werden.


\chapter{Produkteinsatz}
\section{Produktvision}
Das Ziel des Projekts besteht darin, eine Desktopanwendung zu entwickeln, die eine 3D-Ansicht der Welt mit Hilfe freier Daten aus dem OpenStreetMap-Projekt bietet. Die grafische Benutzeroberfläche zeigt dafür eine Weltkugel, die frei gedreht und gezoomt werden kann. 

Als Oberflächentextur kann dafür anfangs auf die freien Satellitenbilder der
NASA zurückgegriffen werden, beim hineinzoomen in die Karte wird dann auf eine Kartenansicht von OpenStreetMap gewechselt.

Die einfach zu bedienende Oberfläche bietet eine simple Steuerung via Maus und Tastatur und erlaubt zudem in verschiedene Einstellungen die Darstellung der Welt zu beeinflussen. 


\section{Anwendungsbereich}



\section{Zielgruppe}
Primäre Zielgruppe des Systems sind Privatpersonen (Jugendliche sowie erwachsene Personen), die eine andere Art der Kartendarstellung als die typischen Onlinekarten bevorzugen.

Eine weitere Zielgruppe sollen wissbegierige Kinder darstellen. Voraussetzung ist lediglich der geübte Umgang mit der Maus und/ oder Tastatur. (Eingeschränkte Features)

\section{Betriebsbedingungen}
Lebensdauer, Ausfallsicherheit, Beaufsichtigung(Wartung)

\begin{itemize}
\item Bestehende dauerhafte Internetverbindung zum Laden des Kartenmaterials.
\item Nach  der  Abschlusspräsentation  werden  von  uns  keine  weiteren 
Veränderungen vorgenommen. Es erfolgt keine Wartung durch uns.
\end{itemize} 

\section{Sicherheit, Datenschutz, gesetzliche Vorgaben}
Verwendung von OpenStreetMap, Jogl usw. ist lizenzfrei...


\chapter{Umgebung}
\section{Software}
Betriebssystem: Windows 7, Windows 8, Linux ...
Jogl, OpenStreetMap, Java 7
soll Plattform-unabhängig sein

\section{Hardware}
Aktueller Standard-PC
Breitbandanbindung, mindestens 1MBit/s
2-4GB RAM

Grafikkarte usw.

Infrastruktur: 1 Server mit 1x Intel Celeron 1200 MHz, 256 MB RAM, 2x40 
GB als RAID 0/1, 100MB Speicherplatz für Daten im Zusammenhang mit 
diesem Projekt 

Verfügbarkeit:  Der  Server  steht  in  den  Geschäftsräumen  von  Frau  Inhalt 
und ist mit ca. 100MBit/s an das Netzwerk Ihres Providers über Richtfunk 
angebunden. 

Internet-Anbindung: je nach Auslastung des Providers, ca. 10MBit/s. 


\section{Orgware}
User: Jugendliche/Erwachsene
User: Kinder
Schnittstellen: Es  werden  diverse  externe  Quellen genutzt, darunter OpenStreetMap Kartenmaterial, die im laufenden Projekt noch konkretisiert werden.


\chapter{Zielsetzungen}
[WIE soll WAS erreicht werden?]
Aufgabenstellung / Zielsetzung

\section{Musskriterien}
\section{Wunschkriterien}
\section{Abgrenzungskriterien}
Kein Login, Authentifizierung, Kennwort-Abfragen, keine Datenspeicherung...



\chapter{Produkt-Funktionen}

\chapter{Produktdaten}

\chapter{Produkt-Leistungen}

\chapter{Benutzungsoberfläche}
1. Screenshot GUI für Jugendliche/Erwachsene
2. Screenshot GUI für Kinder

\chapter{Qualitäts-Zielbestimmung}

\chapter{Globale Testszenarien / Testfälle}

\chapter{Entwicklungsumgebung}
\section{Software}
\begin{itemize}
\item Die Entwicklung erfolgt mittels Eclipse. Eine Installation auf dem Rechner des Auftraggebers ist nicht erforderlich. (Runnable Jar)
\end{itemize}

\section{Hardware}

\section{Orgware}
\begin{itemize}
\item Gruppen-Kommunikation hauptsächlich per Email.
\item Regelmäßige Information des Auftraggebers über die Entwicklungsprozesse/ Phasen/ Ergebnisse.
\end{itemize}




\chapter{Anhang}
\section{Glossar}











\end {document}
